%Stand 18. Juli 29026
\subsection{Anisotropie}
\label{KapitelTheorieAnisotropie}
In Kapitel \ref{KapitelTheorieGrundlagen} wurde bereits erwähnt,
dass ein Semivariogramm \textit{anisotrop} genannt wird, falls 
seine Form nicht nur vom Betrag des 
Abstandsvektors $\Vert s_1 - s_2\Vert$, sondern zusätzlich auch 
von dessen Richtung abhängt.\\
Da sich die mathematische Behandlung einfacher gestaltet, 
ist die Annahme der Isotropie ein sinnvoller Startpunkt, die in realen Anwendungen aber oftmals 
verletzt wird.\\
Gründe hierfür sind beispielsweise
\begin{itemize}[topsep=2pt,parsep=1pt,itemsep=0.5pt]
\item[$\circ$] geologische Strukturen (Hügel, Flussläufe),
\item[$\circ$] bauliche Gegebenheiten (Häuser, Straßen),
\item[$\circ$] unterschiedliche Vegetation (Wälder, Ödland) oder
\item[$\circ$] meteorologische Einflüsse (Windströmungen, Regenwahrscheinlichkeit).
\end{itemize}
\vspace*{1mm}

Man unterscheidet zwischen geometrischer und zonaler Anisotropie 
(vgl.\,Abb.\,(\ref{Abb_Anisotropie})).

\begin{figure}[H] 
\centering
\includegraphics[scale = 0.5]
{Abb_Anisotropie.jpg} 
\caption{\small{Illustration der geometrischen Anisotropie (links) 
mit den beiden unterschiedlichen Werten des Sill (\textcolor{red}{\rule[1mm]{3mm}{0.5mm}}) $a_1$, $a_2$ und der zonalen Anisotropie (rechts)}}
\label{Abb_Anisotropie}
\end{figure} 

$\circ$ \underline{Geometrische Anisotropie}\\
Die geometrische Anisotropie wird dadurch charakterisiert, dass richtungsabhängige 
Variogramme des Zufallsfelds dieselbe Höhe des Sill aufweisen, der Sill aber bei
unterschiedlichen Reichweiten (Ranges) erreicht wird.

Die Korrelation nimmt also in verschiedenen Richtungen unterschiedlich schnell ab.\\

$\circ$ \underline{Zonale Anisotropie}\\
Bei zonaler Anisotropie besitzen die richtungsabhängigen Variogramme
unterschiedliche Sill-Werte, während die Reichweite nahezu identisch bleibt. 
Zonale Anisotropie tritt in Reinform bei realen Phänomenen eher selten 
auf, sondern meist in Kombination mit geometrischer Anisotropie 
(hybride Anisotropie). Anders als die 
geometrische Anisotropie lässt sich die zonale Anisotropie 
\textit{nicht} auf einfache Weise in eine isotrope Struktur 
überführen (vgl.\,hierzu z.\,B.\,\cite[Kap.\,2.5.2]{Chiles1999}, 
\cite[Kap.\,6]{Schmid2025}). \\ 

Um eine mögliche Anisiotropie aufzudecken, wird das Variogramm 
für unterschiedliche  Richtungen berechnet. 
Das Vorgehen zur Schätzung dieser 
\underline{direktionalen Semivarianz} ist der Schätzung 
der empirischen Semivarianz (vgl.\,
Kapitel \ref{KapitelTheorieVariogrammSchätzung}) sehr 
ähnlich. Lediglich findet die Einteilung in 
die Distanzklassen (\textit{Lags}) hier nur für die
 Beobachtungen statt, die sich in einer bestimmten
 Winkelbereich $\vartheta \pm \Delta \vartheta$ befinden.\\
Das genaue Vorgehen ist detaillierter z.\,B.\,in 
\cite[S.\,72ff]{Webster2007}, \cite[S.\,54f]{Montero2015} 
beschrieben. \\

Beispielhaft zeigt 
Abb.\,(\ref{Abb_TheoAnisotropie_Beispiel_XY_Variogramme}) die
Variogramme in X- bzw.\,Y-Richtung des simulierten 
Zufallsfelds aus Abb.\,(\ref{Abb_TheoAnisotropie_BeispielGRF}).

\begin{figure}[H] 
\centering
\includegraphics[scale = 0.48]
{TheoAnisotropie_Beispiel_XY_Variogramme.png} 
\caption{\small{Variogramm in X-Richtung (links) und Y-Richtung
 (rechts) eines geometrisch anisotropen Zufallsfelds}}
\label{Abb_TheoAnisotropie_Beispiel_XY_Variogramme}
\end{figure} 

\begin{figure}[H] 
\centering
\includegraphics[scale = 0.46]
{TheoAnisotropie_BeispielGRF.png} 
\caption{\small{Stichprobe eines geometrisch anisotropen 
Gauß'schen Zufallsfelds ($\lambda=6$)}}
\label{Abb_TheoAnisotropie_BeispielGRF}
\end{figure} 

Liefert das richtungsabhängige Variogramm Hinweise auf eine
geometrische Anisotropie, so kann das zugrundeliegende
Zufallsfeld möglicherweise mittels einer Koordinatentransformation in ein
isotropes Zufallsfeld überführt werden. 
Die für die Umrechnung benötigten Größen sind aus 
Abb.\,(\ref{Abb_Anisotropie_Ellipse}) ersichtlich. 

\begin{figure}[H] 
\centering
\includegraphics[scale = 0.45]
{Abb_Anisotropie_Ellipse.jpg} 
\caption{\small{Skizze der für die Transformation der geometrischen
 Anisotropie verwendeten Größen $a_1$, $a_2$ und $\beta$}}
\label{Abb_Anisotropie_Ellipse}
\end{figure} 

Die Orientierung der Anisotropie $\beta$ wird durch 
den Winkel zwischen der Hauptachse der Ellipse und der
$y$-Achse bestimmt. 
Das Verhältnis der beiden Reichweiten

\begin{equation*}
\lambda =
\frac{a_2}{a_1},\ \lambda > 1,
\end{equation*}
mit 
\begin{itemize}[topsep=2pt,parsep=1pt,itemsep=0.5pt]
\item[$\circ$] der kleineren Reichweite $a_1$ und
\item[$\circ$] der größere Reichweite $a_2$,
\end{itemize}
wird \underline{Anisotropieverhältnis} genannt und ist ein Maß 
für die Richtungsabhängigkeit:
Je größer $\lambda$, desto ausgeprägter ist diese.\\

Zur Modellierung geometrischer Anisotropie wird das 
Koordinatensystem so transformiert, dass die Ellipse in einen Kreis
 überführt wird. Die Zielsetzung der Transformation besteht darin,
die Anisotropie vollständig in die Koordinaten zu verlagern,
sodass anschließend ein isotropes Kovarianz-Modell verwendet werden 
kann.\\

Dies geschieht in zwei Schritten:
\begin{itemize}
\item[ 1)] \underline{Rotation} \\
Zunächst werden die Koordinaten um den Winkel $\beta$
rotiert.
\begin{equation*}
\begin{pmatrix}
x'\\
y'
\end{pmatrix}
=
\begin{pmatrix}
\cos(\beta) & -\sin(\beta)\\
\sin(\beta) & \cos(\beta)
\end{pmatrix}
\begin{pmatrix}
x\\
y
\end{pmatrix}.
\end{equation*}

\item[ 2) ] \underline{Skalierung}\\
Anschließend wird entlang der kurzen Ellipsenachse skaliert:
\begin{equation*}
\begin{pmatrix}
\tilde{x} \\
\tilde{y}
\end{pmatrix}
=
\begin{pmatrix}
1&0\\
0&\lambda
\end{pmatrix}
\begin{pmatrix}
x'\\
y'
\end{pmatrix}.
\end{equation*}
\end{itemize}

Fasst man beide Schritte zusammen, so erhält man als
Transformationsvorschrift

\begin{equation}
\begin{aligned}
\tilde{x}
&=
x\cos(\beta) - y\sin(\beta)\\
\tilde{y}
&=
\lambda x\sin(\beta) +\lambda y\cos(\beta)
\end{aligned}
\end{equation}

wodurch die Ellipse in einen Kreis überführt wird.\\

Im transformierten Koordinatensystem kann anschließend 
ein gewöhnliches isotropes Variogrammmodell verwendet werden.
Bei den im Kapitel \ref{KapitelTheorieKovarianzFunktionen}
eingeführten Kovarianz- und Korrelationsmodellen 
wird Anisotropie dementsprechend durch das Hinzufügen der zwei 
oben eingeführten Parameter modelliert: dem Anisotropieverhältnis
$\lambda$ und dem Winkel $\beta$ der Orientierung. \\

Abschließend sei angemerkt, dass ein richtungsabhängiges 
Variogramm nicht \textit{zwangsläufig} durch Anisotropie verursacht 
sein muss. Auch Verletzungen der Stationaritätsannahme können zu 
einem ähnlichen Verhalten führen (vgl.\,z.\,B.\,\cite[Kap.\,3.7]{Montero2015},
\cite[Kap.\,5]{Kitanidis1997}).